%% Section 4: Intersection theory and height inequality on the total space

We keep the notation of $\mathsection$\ref{SectionSettingUpHtIneq}. 
So we have a closed immersion $\cA\rightarrow \IP_k^n\times
S$ over $S$ satisfying the properties stated near the beginning 
of $\mathsection$\ref{SectionSettingUpHtIneq}. Moreover, $S$ is a
Zariski open subset of an irreducible projective variety $\overline
S\subset\IP_k^m$. 
We assume in addition $k = \IQbar$. Let $X$ be a closed irreducible subvariety of $\cA$ of dimension $d$ defined over $\IQbar$, such that $\pi|_X \colon X\rightarrow S$ is dominant. Let $\omega$ be the Betti form on $\cA$ as defined in Proposition~\ref{PropBettiForm}.
\begin{proposition}\label{PropAuxHtIneq}
We keep the notation from above and
suppose that $\an{X}$ contains a smooth point at which
$\omega|_{\sman{X}}^{\wedge d} >0$. Then there exists a constant $c_1
> 0$ satisfying the following property.
Let $N \in \IN$ be a power of $2$, there exist a Zariski open dense subset $U_N$ of $X$ defined over $\IQbar$ and a constant $c_2(N)$ such that
\[
h([N]P) \ge c_1 N^2 h(\pi(P)) - c_2(N)\qquad \text{for all }P \in U_N(\IQbar).
\]
\end{proposition}

The goal of this section is to prove Proposition~\ref{PropAuxHtIneq}. 
The key idea is to apply a theorem of Siu \cite[Theorem~2.2.15]{PosAlgGeom}. Let us briefly explain the main points before moving on to the proof.

Let $X$ be as in Proposition~\ref{PropAuxHtIneq}, and let $P \in
X(\IQbar)$. For each $N \in \IN$, we work with $X_N \subseteq
\cA\times_S \cA$, the graph of $[N] \colon \cA \rightarrow \cA$, and
its Zariski closure $\overline{X_N}$ in
$\overline{\cA}\times_{\overline{S}}\overline{\cA} \subseteq \IP_k^n
\times \IP_k^n \times \IP_k^m$. The point $P$ gives rise to a point
$P_N \in \overline{X_N}$; see $\mathsection$\ref{ss:graph}. Consider
the line bundles $\cF = \cO(0,1,1)|_{\overline{X_N}}$ and $\cM =
\cO(0,0,1)|_{\overline{X_N}}$. Choosing representatives as in
last paragraph of
$\mathsection$\ref{ss:graph}  
 our height inequality in Proposition~\ref{PropAuxHtIneq} is equivalent to
\[
h_{\overline{X_N},\cF}(P_N)\ge c_1 N^2 h_{\overline{X_N},\cM}(P_N) - c_2'(N)
\]
for some $c_2'(N)$ independent of $P$ (which may be different from
$c_2(N)$). By the Height Machine  it suffices to find positive
integers $p$ and $q$, independent of $N$, such that $\cF^{\otimes q}
\otimes \cM^{\otimes -p N^2}$ is a big line bundle on
$\overline{X_N}$; we can then take $c_1=p/q$. 

Both $\cF$ and $\cM$ are nef line bundles. Thus a criterion of bigness
by Siu \cite[Theorem~2.2.15]{PosAlgGeom}, states that $\cF^{\otimes q}
\otimes \cM^{\otimes -p N^2}$ is big if $(\cF^{\cdot d}) > d c_1
(\cM^{\otimes N^2} \cdot \cF^{\cdot (d-1)})$.
Note that $(\cM^{\otimes N^2} \cdot \cF^{\cdot (d-1)}) = N^2 (\cM \cdot \cF^{\cdot (d-1)})$ by multi-linearity of intersection numbers. 
Thus our task 
becomes comparing two intersection numbers. Our application continues
to work if the numerical
factor $d = \dim X$ is replaced by any
positive  factor that depends only on the dimension.
So it remains to prove an appropriate lower bound for $(\cF^{\cdot d})$ and an appropriate upper bound for $(\cM \cdot \cF^{\cdot (d-1)})$.

The proof of Proposition~\ref{PropAuxHtIneq} will be organized as follows in this section. We first prove the appropriate lower bound for $(\cF^{\cdot d})$ in Proposition~\ref{prop:intersectionlb}. This is where we use the hypothesis that $\omega|_{\sman{X}}^{\wedge d} >0$ at some smooth point of $\an{X}$. Next we prove the appropriate lower bound for $(\cM \cdot \cF^{\cdot (d-1)})$ in Proposition~\ref{prop:intersectionub}. At this step the assumption of $N$ being a power of $2$ is used. Then we finish the proof of Proposition~\ref{PropAuxHtIneq} by applying Siu's theorem in $\mathsection$\ref{SubsectionProofOfAUxHtIneq}.




\subsection{Bounding an intersection number from below}
Let $X$ be as in Proposition~\ref{PropAuxHtIneq}. 
For each $N \in \IN$, let $\overline{X_N} \subseteq \IP_k^n \times
\IP_k^n \times \IP_k^m$ be as in $\mathsection$\ref{ss:graph}.
In particular, $\dim \overline{X_N} = d$.
%Without loss of generality we may assume $k=\IC$. 
 Let $\cF =
\cO(0,1,1)|_{\overline{X_N}}$ be as in \eqref{def:cF}. The top
self-intersection of $\cF$ on $\overline{X_N}$ is bounded from below
in the following proposition.
To prove it, we may replace $X$ by its base change to $\IC$. 

\begin{proposition}
  \label{prop:intersectionlb}
  Suppose $\an{X}$ contains a smooth point at which
  $\omega|_{\sman{X}}^{\wedge d} >0$. Then there exists a constant $\kappa
  > 0$, independent of $N$, such that $(\cF^{\cdot d}) \ge \kappa N^{2d}$ for all
  $N\in\IN$.   
\end{proposition}
\begin{proof}
We fix a point  $P_0 \in \sman{X}$ at which
$\omega|_{\sman{X}}^{\wedge d}$ is positive and let $s_0=\pi(P_0),\Delta,\vartheta,\theta,$
and $K$  be as in $\mathsection$\ref{SubsectionSetUpHtIneqBetti}, see 
Remark~\ref{rmk:nondeg}. In particular $\vartheta(P_0) = \theta\circ \pi(P_0) = 1$. 
We extend $\vartheta$ to a smooth function on  $\an{(\IP^m_\IC)}$ by
setting it $0$ outside of the compact set $K\subset \an{S}$. This
extends $\theta = \vartheta\circ\pi$ to all of
$\an{(\IP^n_\IC\times\IP^m_\IC)}$. 
  %%   As stated above, $\int_{\sman{X}} \omega^{\wedge d}$ is not
%% well-defined if
%% $\omega$ does not have compact support. However, $\omega_c$ as in
%% (\ref{def:omegac}) does  have
%% compact support on $\cA$. 

Let $\alpha$ be the pull-back of the Fubini--Study form 
 under the analytification of the Segre morphism
$\IP_{\IC}^n\times\IP_{\IC}^m \rightarrow \IP_{\IC}^{(n+1)(m+1)-1}$.
We replace $\alpha$ by its restriction to $\an{\overline\cA}$. 
Thus $\alpha$
represents the Chern class of
$\cO(1,1) \in \mathrm{Pic}(\IP^n_{\IC}\times\IP_{\IC}^m)$ restricted
 to $\an{\overline\cA}$, using common notation.

%% Recall that  $\alpha$ represents the Chern class of an ample line
%% bundle on ${\cA}$. So it is
Note that $\alpha$ is strictly positive on all of
$\an{\cA}$.  Since $\Delta$ is relatively compact we can find a constant $C>0$
with
\begin{equation}\label{EqConditionOnCapitalC}
 C \alpha|_{\cA_{\Delta}} - \omega|_{\cA_{\Delta}} \ge 0.
\end{equation}
As the smooth and non-negative function $\theta = \vartheta\circ\pi$ on $\an{\cA}$ has support in
$\pi^{-1}(K)\subset \pi^{-1}(\Delta)=\cA_{\Delta}$ we have
\begin{equation*}
  C \theta\alpha - \theta\omega \ge 0.
\end{equation*}

We pull this $(1,1)$-form back under the holomorphic map
$[N]\colon \an{\cA}\rightarrow\an{\cA}$ and get
\begin{equation}
\label{eq:pbsemipositive}
C  [N]^* (\theta\alpha) - N^2\theta\omega =  C [N]^*(\theta\alpha) - [N]^*(\theta\omega)\ge 0
\end{equation}
where we used $[N]^*\omega = N^2\omega$ and $[N]^* \theta = \theta$; 
the former is a property of the Betti form, see
Proposition~\ref{PropBettiForm}(ii) and
the latter holds as $\theta$ is the pull-back from the base of $\vartheta$.

We define
\begin{equation*}
  \beta = C  [N]^*(\theta\alpha) - N^2\theta\omega,
\end{equation*}
which is a  $(1,1)$-form on $\an{\cA}$.
It is semi-positive by (\ref{eq:pbsemipositive}). 
The support of $\theta$ is
contained in $\pi^{-1}(K)$, which we have identified as compact at the
end of $\mathsection$\ref{SubsectionSetUpHtIneqBetti}.
So
$C[N]^*(\theta\alpha)$ and $N^2\theta\omega$ have compact support on $\an{\cA}$. % Thus we are allowed to integrate.

%\begin{lemma}
We claim that  $\int_{\sman{X}} (C[N]^*(\theta\alpha))^{\wedge d} \ge
\int_{\sman{X}} (N^2 \theta\omega)^{\wedge d}$.

First observe that both integrals are well-defined as both
$[N]^*(\theta\alpha)$ and $N^2\theta\omega$ have compact support on
$\an{\cA}$; this follows from work of Lelong~\cite{Lelong} which we
use freely below. A textbook proof can be found in
\cite[Theorem~11.21]{voisin:hodge1} and
\cite[$\mathsection$III.2.B]{Demailly}. 
%\end{lemma}
%\begin{proof}
To prove the inequality let us  write $\beta = \gamma - \delta$ with $\gamma = C[N]^*(\theta
\alpha)$ and $\delta = N^2\theta\omega$.
% all of compact support, [PH]removed, as redundant
Then
\begin{equation}
\label{eq:diffintegral}  
  \int_{\sman{X}} \gamma^{\wedge d} - \int_{\sman{X}}\delta^{\wedge d}
  = 
\int_{\sman{X}} (\delta+\beta)^{\wedge d} - \int_{\sman{X}} \delta^{\wedge
    d}
= \sum_{i=0}^{d-1} {d \choose i} \int_{\sman{X}} \delta^{\wedge
  i}\wedge \beta^{\wedge (d-i)}
\end{equation}
as the exterior product is commutative on even degree forms. 
We know that $\beta \ge 0$ on $\an{\cA}$ and it is also crucial that
$\delta \ge 0$ on $\an{\cA}$, the latter
follows from $\omega \ge 0$, property (i) of
Proposition~\ref{PropBettiForm},
and from $\theta\ge 0$.
Then $\delta^{\wedge i}\wedge \beta^{\wedge (d-i)}$ is semi-positive
on $\an{\cA}$; see \cite[Proposition~III.1.11]{Demailly}\footnote{As our convention  is somewhat different from Demailly's, let us explain how to apply  \cite[Proposition~III.1.11]{Demailly}. Our definition of semi-positive $(1,1)$-form coincides with that of \textit{positive $(1,1)$-form} of \cite[Chapter~III]{Demailly} by Corollary~1.7 of \textit{loc.cit.}, and thus are precisely the strongly positive $(1,1)$-forms of \cite[Chapter~III]{Demailly} by Corollary~1.9 of \textit{loc.cit.} Therefore we can apply the cited proposition.}. Thus the right-hand side of (\ref{eq:diffintegral}) is
 non-negative, see \cite[Theorem~III.2.7]{Demailly}, and our claim is settled.
%\end{proof}

The claim implies
\begin{equation}
\label{eq:N2dgrowth}
 C^d \int_{\sman{X}}  [N]^*(\theta\alpha)^{\wedge d} \ge
\kappa' N^{2d} \quad\text{where}\quad \kappa' = \int_{\sman{X}} (\theta\omega)^{\wedge d}.
\end{equation}
We have $\kappa' > 0$. Indeed, $(\theta\omega)^{\wedge d}$ is semi-positive on $\an{\cA}$ because $\omega \ge 0$ (Proposition~\ref{PropBettiForm}(i)) and $\theta \ge 0$ (by construction). But $\omega|_{\sman{X}}^{\wedge d}$ is positive at $P_0 \in \sman{X}$ by choice of $P_0$ and $\theta\circ \pi(P_0) = 1$ by choice of $\theta$. So $(\theta\omega)|_{\sman{X}}^{\wedge d}$ is positive at $P_0 \in \sman{X}$. Thus $\kappa' > 0$.

Next we want to relate the integral on the left in (\ref{eq:N2dgrowth}) with an
intersection number. First we recall that $[N]$ is given in terms of
the graph construction, \textit{cf.} (\ref{eq:relatetoN}). So we may
rewrite
\begin{equation}
\label{eq:chainbegins}
  \int_{\sman{X}} [N]^*(\theta\alpha)^{\wedge d}  =
  \int_{\sman{X}} (\rho_2|_{\Gamma_N} \circ
  \rho_1|_{\Gamma_N}^{-1})^*(\theta\alpha)^{\wedge d}=
  \int_{\sman{X}} (\rho_1|_{\Gamma_N}^{-1})^*\rho_2|_{\Gamma_N}^*(\theta\alpha)^{\wedge d},
\end{equation}
here $\Gamma_N,\rho_1,$ and $\rho_2$ are as defined in $\mathsection$\ref{ss:graph}. 

Because $\rho_1|_{\an{\Gamma}_N} \colon \an{\Gamma}_N \rightarrow \an{\cA}$ is biholomorphic we
can change coordinates and integrate over $X_N$, which is a complex
analytic subset of the graph $\Gamma_N$, itself a complex manifold.
More precisely, we have
\begin{equation}
\label{eq:chain2} 
  \int_{\sman{X}} (\rho_1|_{\Gamma_N}^{-1})^*\rho_2|_{\Gamma_N}^*(\theta\alpha)^{\wedge d}
=\int_{\sman{X}_N}  \rho_2|_{\Gamma_N}^*(\theta\alpha)^{\wedge d}.
\end{equation} 


Recall that $\alpha$ is the restriction to $\an{\overline{\cA}}$ 
of a  $(1,1)$-form on  
$\an{(\IP^n_{\IC}\times\IP^m_{\IC})}$.
Moreover,
$\rho_2$ is also defined on all of $\IP^n_{\IC} \times\IP^n_{\IC}\times \IP^m_{\IC}$. So ${\rho_2}|_{\Gamma_N}^*(\theta\alpha)$ is the restriction to
$\Gamma_N$ of a $(1,1)$-form defined on $\an{(\IP^n_{\IC} \times\IP^n_{\IC}\times \IP^m_{\IC})}$. 
Observe that $\sman{X}_N\subset \an{\overline{X_N}}$ and the difference has
dimension strictly less than $d=\dim X_N$.
This justifies% the first inequality in 
\begin{equation}
\label{eq:chain3}   
  \int_{\sman{X}_N} \rho_2|_{\Gamma_N}^*(\theta\alpha)^{\wedge d}
  =\int_{\an{\overline{X_N}}} \rho_2^*(\theta\alpha)^{\wedge d}
%  \le \int_{\an{\overline{X_N}}} \rho_2^*\alpha^{\wedge d}
\end{equation}
where we take $\an{\overline{X_N}}$ as a complex analytic subset of
the analytification of $\IP_{\IC}^n\times\IP_{\IC}^n\times \IP_{\IC}^m$ and $\rho_2^*(\theta\alpha)$ as a
$(1,1)$-form on this ambient space.
Now $\theta$ takes values in $[0,1]$ and so
\begin{equation}
\label{eq:chain4}
\int_{\an{\overline{X_N}}} \rho_2^*(\theta\alpha)^{\wedge d}\le
%  \int_{\an{\overline{X_N}}} \theta^d(\rho_2^*\alpha)^{\wedge d}\le
  \int_{\an{\overline{X_N}}} (\rho_2^*\alpha)^{\wedge d}.
\end{equation}


The pull-back $\rho_2^*\alpha$ represents
$\rho_2^* \cO(1,1) \in
\mathrm{Pic}(\IP_{\IC}^n\times\IP_{\IC}^n\times\IP_{\IC}^m)$ in the
Picard group and
has compact support as $\an{(\IP_{\IC}^n\times\IP_{\IC}^n\times \IP_{\IC}^m)}$ is
compact.
But integration coincides with the intersection pairing in the compact
case; see \cite[Theorem~11.21]{voisin:hodge1}. In particular, we have
\begin{equation}
\label{eq:chain5} 
\int_{\an{\overline{X_N}}} (\rho_2^*\alpha)^{\wedge d} = (\rho_2^*
\cO(1,1)^{\cdot d} [\overline {X_N}]) %= \left(\cO(1,1,0)^{\cdot d}[\overline {X_N}]\right)
\end{equation}
where the intersection takes place in
$\IP_{\IC}^n\times\IP_{\IC}^n\times\IP_{\IC}^m$. We recall \eqref{def:cF}
and apply  the projection formula to obtain 
\begin{equation}
\label{eq:chain6}  
  (\rho_2^*
\cO(1,1)^{\cdot d}  [\overline {X_N}]) = (\cO(0,1,1)^{\cdot d}  [\overline {X_N}]) =  (\cF^{\cdot d}).
\end{equation}
%where we use $(\cF^{\cdot d})$ to denote  the self-intersection number on $\overline{X_N}$. 

The (in)equalities  (\ref{eq:chainbegins}),
(\ref{eq:chain2}), (\ref{eq:chain3}),
(\ref{eq:chain4}), (\ref{eq:chain5}), and (\ref{eq:chain6}) yield
\begin{equation*}
  \int_{\sman{X}} [N]^*(\theta\alpha)^{\wedge d} \le (\cF^{\cdot d}). 
\end{equation*}
We recall the  lower bound (\ref{eq:N2dgrowth}) % and the definition
% (\ref{def:cF})
to  obtain
$(\cF^{\cdot d}) \ge (\kappa'/C^d) N^{2d}$ where $C$ comes from
\eqref{EqConditionOnCapitalC} and $\kappa' > 0$ comes from
\eqref{eq:N2dgrowth}.
The proposition follows with $\kappa = \kappa'/C^d$. 
\end{proof}


\subsection{Bounding an intersection number from above}

We keep the notation from the last subsection with $k=\IQbar$. 
So  $X$ is as above Proposition~\ref{PropAuxHtIneq} with $\dim X=d$
For each $N \in \IN$, let $\overline{X_N} \subseteq \IP_k^n \times
\IP_k^n \times \IP_k^m$ be the graph construction as in $\mathsection$\ref{ss:graph}. In
particular, $\dim \overline{X_N} = d$.  Here we need $\cF =
\cO(0,1,1)|_{\overline{X_N}}$ as defined in (\ref{def:cF}) and also
 $\cM = \cO(0,0,1)|_{\overline{X_N}}$ as defined in \eqref{def:cM}.

\begin{proposition}
  \label{prop:intersectionub}
Assume $d \ge 1$.  There exists a constant $c>0$ depending on the data introduced above
  with the following property. Say $N\ge 1$ is a power of $2$, then
  \begin{equation*}
    (\cM\cdot \cF^{\cdot(d-1)}) \le c N^{2(d-1)}.
  \end{equation*}
% where  $\cM =\rho_0|_{\overline{X_N}}^*\cO(1)$.
\end{proposition}

Let us make some preliminary remarks before the proof. A similar upper
bound for the intersection number was derived by the third-named
author in \cite{habegger:imrn,Hab:Special} using Philippon's
version~\cite{Philippon} of B\'ezout's Theorem for multiprojective
space. The approach here is similar but does not
 refer to Philippon's result. Rather, we rely on the following well-known
positivity property of the intersection theory of multiprojective
space: any effective Weil divisor on a multiprojective space is
nef. This approach was motivated by
K\"uhne's~\cite{kuehne:semiabelianbhc} work on semiabelian varieties.
% We refer to Appendix~\ref{app:positivity} for
%  the intersection theory we need.

\begin{proof}[Proof of Proposition \ref{prop:intersectionub}] Recall that $[2] \colon
A \rightarrow A$  is the  multiplication-by-$2$ morphism on $A$. For the symmetric and ample line bundle
$L$ on $A$, we have $[2]^*L \cong L^{\otimes 4}$.
Recall that $A$ is projectively normal in $\IP_{k(S)}^n$. 
By a result of Serre, \cite[Corollaire~2, Appendix~II]{Ast6970}, 
the morphism $[2]$ is represented 
 by homogeneous
polynomials $f_0,\ldots,f_n$ in the $n+1$ projective coordinates of $\IP^n$ of degree $4$, with  coefficients in $k(S)$ and with no common
zeros in $A$. 

Recall that the family $\cA$ is embedded in
$\IP_k^n\times S \subset\IP_k^n\times\IP_k^m$. We can spread out the
$f_0,\ldots,f_n$. More precisely, there exist a Zariski closed, proper
subset $Z\subsetneq \cA$ and polynomials $f_0,\ldots,f_n \in
k[\mathbf{X},\mathbf{S}]$ that are bihomogeneous of degree $(4,D')$ in
 the $(n+1)$-tuple of projective coordinates $\mathbf{X}$ of
$\IP_k^n$
and the $(m+1)$-tuple of projective coordinates $\mathbf{S}$ of
$\IP_k^m$,
 with the following properties:
\begin{enumerate}
 \item [(i)] the polynomials $f_0,\ldots,f_n$ have no common zeros on
 $(\cA\setminus Z)(k)$, and 
 \item[(ii)]  if   $(P,s)\in(\cA\setminus Z)(k)$, then $[2](P,s) = \bigl([f_0(P,s):\cdots : f_n(P,s)],s\bigr)$. 
\end{enumerate}
Moreover, as $f_0,\ldots,f_n$ have no common zero on the generic
fiber, we may assume that $\pi(Z)$ is Zariski closed and proper in
$S$. So we may assume that $Z = \pi^{-1}(\pi(Z))\subsetneq \cA$ and in particular,
$[2]$ maps $\cA\setminus Z$ to itself.

The  $4$ in the bidegree $(4,D')$ comes from $2^2=4$. The degree
$D'$ with respect to the base coordinates $\mathbf{S}$ is more mysterious. However,
by successively iterating we will get it under control. 

%% \begin{remark}
%% This degree may be a source of ineffectivity! But it is probably
%% computable for $\mathcal{A}_g$.   
%% \end{remark}

For each integer $l\ge 1$ we require polynomials
$f^{(l)}_0,\ldots,f^{(l)}_n$ to describe multiplication-by-$2^l$,
\textit{cf.} \cite[$\mathsection$9]{GaoHab}. In order to obtain information on the
degree with respect to $\mathbf{S}$ we construct them by iterating
the $f^{(1)}_0=f_0,\ldots,f^{(1)}_n=f_n$.
For all $i\in \{0,\ldots,n\}$  we set
\begin{equation*}
%  f^{(l+1)}_i((S_j),(X_j)) =
%  f\left((S_j)_j,\left(f^{(l)}_0((S_j)_j,(X_j)_j),\ldots,f^{(l)}_n((S_j)_j,(X_j)_j)\right)\right)
  f^{(l+1)}_i(\mathbf{X},\mathbf{S}) = f_i\left(\bigl(f^{(l)}_0(\mathbf{X},\mathbf{S}),\ldots,f^{(l)}_n(\mathbf{X},\mathbf{S})\bigr),\mathbf{S}\right)
\end{equation*}
 for all $i$; it is bihomogeneous in
$\mathbf{X}$ and $\mathbf{S}$. 
So for all $l\ge 1$ 
\begin{enumerate}
 \item [(i)] the polynomials $f^{(l)}_0,\ldots,f^{(l)}_n$ have no common zeros on
 $(\cA\setminus Z)(k)$, and 
 \item[(ii)]  if   $(P,s)\in(\cA\setminus Z)(k)$, then $[2^l](P,s) = \bigl([f^{(l)}_0(P,s):\cdots : f^{(l)}_n(P,s)],s\bigr)$. 
\end{enumerate}


If for all $i$ the polynomials $f^{(l)}_i$ are bihomogeneous of degree
$(D_{l},D'_{l})$, then all $f^{(l+1)}_i$ are bihomogeneous of degree
$(4D_{l},D'+4D'_{l})$. Recall that
$(D_1,D'_1)=(4,D')$, thus the recurrence
relations
$$ D_{l+1}=4D_l  \quad\text{and}\quad D'_{l+1} = D' + 4D'_{l} $$
imply
\begin{equation}
  \label{eq:DlprimeDlbound}
  D_l = 4^l \quad\text{and}\quad   D'_l = \frac{4^l-1}{3}D'\le 4^lD'
\end{equation}
for all $l\ge 1$.
Up-to the constant linear factor $D'$ the bidegrees both
grow like $4^l$.


We proceed as follows to cut out the graph $\overline{X_{N}}$ where
$N=2^l$. 
We start out with $\overline X\subset\IP_k^n\times\IP_k^m$. 
As $X$ dominates $S$ but $Z$ does not, there is an $i$ such that 
$f_i^{(l)}$  does not
vanish identically on $X$, without loss of generality we assume $i=0$. 


Then as $i$ varies over $\{1,\ldots,n\}$ we obtain $n$ trihomogeneous
polynomials 
\begin{equation*}
g_i :=  Y_i f_0^{(l)}(\mathbf{X},\mathbf{S}) - Y_0   f_i^{(l)}(\mathbf{X},\mathbf{S})
\end{equation*}
where $Y_0,\ldots,Y_n$ are the projective coordinates on the middle factor of $\IP_k^n\times\IP_k^n\times\IP_k^m$. 
The tridegree of these polynomials is $(D_l,1,D'_l)$. %\Pcomment{There
  %was a mistake here, $(D'_l,1,D_l)$ instead of $(D_l,1,D'_l)$.}
Their zero locus on $\overline X\times\IP_k^n$ has the graph
$\overline{X_N}$ as an irreducible component;
by permuting coordinates we consider $\overline X\times\IP_k^n$ as a
subvariety of $\IP_k^n\times\IP_k^n\times\IP_k^m$.
We will see below that this is a proper component of the said
intersection. 
However, there may be further irreducible components in this
intersection,  some could even have dimension greater than $\dim
\overline{X_N}$.

%Observe that $\dim\overline{X_N} =\dim\overline X = d$. 
This issue is clarified by the positivity result \cite[Corollary
12.2.(a)]{Fulton}. We apply it to the ambient
variety 
$\IP_k^n\times\IP_k^n\times\IP_k^m$, which becomes $X$ in Fulton's
notation; observe that
 the tangent bundle of a product of projective spaces
is generated by its global sections, \textit{cf.}~\cite[Examples 12.2.1.(a) and
(c)]{Fulton}. For $i\in\{1,\ldots,n\}$, the $V_i$ in Fulton's notation
is the zero set of $g_i$, and $V_{n+1}$ is 
$\overline{X} \times \IP_k^n$. So $r=n+1$ and $V_1,\ldots,V_{n+1}$ are
equidimensional. 
Observe that
\begin{alignat*}1
  \sum_{i=1}^{r} \dim V_i  - (r-1) \dim
  \IP_k^n\times\IP_k^n\times\IP_k^m
  &= (2n+m-1)(r-1) + \dim\overline{X}\times\IP_k^n
  - (r-1)(2n+m)
  \\
  &= \dim \overline{X} = \dim \overline{X_N},
\end{alignat*}
% Observe that $\dim \overline{X_N} = \dim X
% = \dim (\overline{X} \times \IP_k^n) - n$,
so, and as announced above, $\overline{X_N}$ is a proper component in the intersection of
$V_1,\ldots,V_n,$ and $\overline{X} \times \IP_k^n$.
By Fulton's \cite[Corollary
12.2.(a)]{Fulton} the  cycle class
attached to the intersection of $\overline X\times\IP_k^n$ with the zero
locus of $g_1,\ldots,g_n$
is represented by a positive cycle on
$\IP_k^n\times\IP_k^n\times\IP_k^m$, one of
whose components is $\overline{X_N}$. As $\cO(0,0,1)$ and $\cO(0,1,1)$
are numerically effective we conclude 
\begin{alignat}1
  \label{eq:inupperboundconclusion1}
  (\cO(0,0,1)  \cO(0,1,1)^{\cdot (d-1)}  [\overline{X_N}])
  &\le
  \left(\cO(0,0,1)  \cO(0,1,1)^{\cdot (d-1)}
  \cO(D_l,1,D'_l)^{\cdot n}  [\overline{X} \times\IP_k^n]\right).
%\\
 % &\le 4^{l(d-1)}{D'}^{d-1}{n+d\choose n}(\cO(1,1)^{\cdot d}[\overline{X}])
\end{alignat}

The cycle $[\overline X\times\IP_k^n]$ is linearly equivalent to 
$\sum_{i+p=n+m-d} a_{ip}H_1^{\cdot i}  H_2^{\cdot p}$, with $H_{1}$ and $H_2$
hyperplane pullbacks of the factors $\IP_k^n\times \IP_k^m\supset
\overline X$, respectively, and with
$a_{ip}$ non-negative integers that depend only on $\overline X$.
Thus the left-hand side of (\ref{eq:inupperboundconclusion1}) is at
most %\Pcomment{There was a mistake here, it read $\cO(D'_l,1,D_l)$
%  instead of $\cO(D_l,1,D'_l)$.}
\begin{alignat*}1
\sum_{i+p = n+m-d} a_{ip}
\left(  \cO(0,0,1)\cO(0,1,1)^{\cdot (d-1)}
\cO(D_l,1,D'_l)^{\cdot n}\cO(1,0,0)^{\cdot i}\cO(0,0,1)^{\cdot p}\right).
\end{alignat*}
We can expand the sum using linearly of intersection
numbers to find that it equals
\begin{alignat*}1
  \sum_{\substack{i+p=n+m-d\\ j'+p'=d-1\\i''+j''+p'' = n}} a_{ip} {d-1 \choose j',p'}
  {n\choose i'',j'',p''}{D_l}^{i''} {D'_l}^{p''} \left(\cO(1,0,0)^{\cdot
    (i+i'')} \cO(0,1,0)^{\cdot (j'+j'')}\cO(0,0,1)^{\cdot(1+p+p'+p'')}\right)
\end{alignat*}
Only terms with $i+i''\le n$ and $j'+j''\le n$ and $1+p+p'+p''\le m$
contribute to the sum.
On the other hand, any term in the sum satisfies
$i+i'' + j'+j'' + 1+p+p'+p'' =2n+m$. 
So we can assume $i+i''=n$ and $j'+j''=n$ and $1+p+p'+p''=m$ in the
sum
which thus simplifies to
\begin{alignat*}1
  \sum_{\substack{i+p=n+m-d, i+i''=n \\ j'+p'=d-1,j'+j''=n \\ i''+j''+p''=n,p+p'+p''=m-1}} a_{ip} {d-1\choose
    j',p'} {n\choose i'',j'',p''} {D_l}^{i''} {D'_l}^{p''}. 
\end{alignat*}
Note $i''+p''  = n-j'' =j'  = d-1-p'\le d-1$. 
We recall (\ref{eq:DlprimeDlbound}) and conclude that the 
 left-hand side of (\ref{eq:inupperboundconclusion1}) is at most
\begin{equation}
  \label{eq:intersectionnumberub2}
  (4^lD')^{d-1}
  \sum_{\substack{i+p=n+m-d\\ j'+p'=d-1\\i''+j''+p'' = n}}
 a_{ip}
 {d-1\choose
    j',p'} {n\choose i'',j'',p''} 
\le
 (4^l D')^{d-1}2^{d-1}3^n \sum_{\substack{i+p=n+m-d}}
 a_{ip}.
\end{equation}
%where $c>0$ depends only on $X$. 

% $\cO(0,1,1) = \cO(0,1,0)\otimes \cO(0,0,1),
% \cO(D'_l,1,D_l) = \cO(1,0,0)^{\otimes D'_l} \otimes\cO(0,1,0)\otimes
% \cO(0,0,1)^{D_l}$.  Any occurrence $\cO(1,0,0)^{\cdot i}$ or
% $\cO(0,1,0)^{\cdot i}$ with $i>n$ or $\cO(0,0,1)^{\cdot i}$ with $i>m$
% does not contribute to the intersection number. 

We recall $N=2^l$ and use the projection formula with the estimates
above to find 
\begin{equation*}
  (\cO(0,0,1)|_{\overline{X_N}} 
  \cO(0,1,1)|_{\overline{X_N}}^{\cdot(d-1)}) \le  c N^{2(d-1)}
\end{equation*}
where $c>0$ depends only on $X$. 
Recall our definition $\cF = \cO(0,1,1)|_{\overline{X_N}}$ and $\cM =
\cO(0,0,1)|_{\overline{X_N}}$. So we get 
$  (\cM\cdot \cF^{\cdot(d-1)}) \le c N^{2(d-1)}$, as desired. 
%Thus we are done by taking $c = {D'}^{d-1}{n+d\choose n}(\cO(1,1)^{\cdot d}[\overline{X}])$.
%
% We recall $N=2^l$ and use the projection formula and find 
% \begin{equation*}
%   (\cO(0,0,1)|_{\overline{X_N}} \cdot 
%   \cO(0,1,1)|_{\overline{X_N}}^{\cdot(d-1)}) \le  N^{2(d-1)}
%       {D'}^{d-1}
%       {n+d\choose n}(\cO(1,1)^{\cdot d}[\overline{X}])
% \end{equation*}
%
% Recall our definition $\cF = \cO(0,1,1)|_{\overline{X_N}}$ and $\cM =
% \cO(0,0,1)|_{\overline{X_N}}$. So we get 
% \begin{equation}
% \label{eq:intersectionup}
%   (\cM\cdot \cF^{\cdot(d-1)}) \le N^{2(d-1)} {D'}^{d-1}{n+d\choose n}(\cO(1,1)^{\cdot d}[\overline{X}]).
% \end{equation}
% Thus we are done by taking $c = {D'}^{d-1}{n+d\choose n}(\cO(1,1)^{\cdot d}[\overline{X}])$.
\end{proof}


\subsection{Proof of Proposition~\ref{PropAuxHtIneq}}\label{SubsectionProofOfAUxHtIneq}
Now let us prove Proposition~\ref{PropAuxHtIneq} by comparing the intersection number inequalities in
Propositions~\ref{prop:intersectionlb} and \ref{prop:intersectionub}.

Let $X$ be of dimension $d$ as in Proposition~\ref{PropAuxHtIneq}. The
case $d = 0$ is trivial. So we assume $d \ge 1$. We may assume $N=2^l$
with $l \in \IN$. Let $\overline{X_N} \subseteq \IP_k^n \times \IP_k^n \times \IP_k^m$ be as in $\mathsection$\ref{ss:graph}. In particular $\dim \overline{X_N} = d$.

Let $\kappa > 0$ be as in Proposition~\ref{prop:intersectionlb}. Then
$(\cF^{\cdot d}) \ge \kappa N^{2d}$. Let $c > 0$ be as in
Proposition~\ref{prop:intersectionub}. Then $(\cM \cdot \cF^{\cdot
(d-1)}) \le c N^{2(d-1)}$.
We have indicated how to obtain  $\kappa$ and $c$  at the end of the proof of each one of the corresponding propositions.

Fix a rational number $c_1$ such that 
\begin{equation}
  \label{EqConstantInHtIneq}
  0<  c_1  cd < \kappa.
\end{equation}
Let $q$ be a multiple of the denominator of $c_1$.
Using the bounds above and linearity of intersection numbers we get
\begin{equation*}
d (\cM^{\otimes qc_1 N^2}\cdot(\cF^{\otimes q})^{\cdot(d-1)}) = d q^d
c_1 N^2 (\cM \cdot
\cF^{\cdot (d-1)}) \le d q^d c_1 N^2 c N^{2(d-1)} < \kappa q^d  N^{2d} \le
((\cF^{\otimes q})^{\cdot d}).
\end{equation*}
Then $\cF^{\otimes q} \otimes \cM^{\otimes -qc_1N^2}$ is a big line
bundle on $\overline{X_N}$ by a theorem of Siu \cite[Theorem
2.2.15]{PosAlgGeom}.  After possibly replacing $q$ by a multiple the line bundle
 $\cF^{\otimes q} \otimes \cM^{\otimes -qc_1N^2}$ admits a
non-zero global section. Say $h_{\overline{X_N},\cF}$ and
$h_{\overline{X_N},\cM}$ are representives of heights on
$\overline{X_N}$ attached by the Height Machine to
$\cF$ and $\cM$, respectively. After canceling $q$ we conclude that
$h_{\overline{X_N},\cF}- c_1 N^2 h_{\overline{X_N},\cM}$ is bounded
from below on a Zariski open and dense subset of $\overline{X_N}$. 
The  image  of this  subset under the projection $\rho_1$
contains  a Zariski open and dense subset $U_N$ of $X$.
It follows from the end of  $\mathsection$\ref{ss:graph}  that
%% $X(\IQbar)\ni P\mapsto
%% q h([N](P))$ represents the height provided by the Height Machine for
%% $(\overline{X_N},\cF^{\otimes q})$ and $P\mapsto qc_1 N^2 h(\pi(P))$ represents
%% the height attached to $(\overline{X_N},\cM^{\otimes qc_1 N^2})$.
%% Thus there exist a Zariski open dense subset
%% %$U'_l$ of $\overline{X_N}$ defined over $\IQbar$ and a constant
%% $U_N$ of $\overline{X_N}$ defined over $\IQbar$ and a constant
there exists $c_2(N)$ such that
\begin{equation*}%\label{EqAuxHtIneqAux}
%% h_{\overline{X_N},\cF}(P_N) \ge 
%% c_1N^2 h_{\overline{X_N},\cM}(P_N) - c_2(l)\quad \text{ for all }P_N \in U'_l(\IQbar)
h([N](P)) \ge 
c_1N^2 h(\pi(P)) - c_2(N)\quad \text{ for all }P \in U_N(\IQbar). \makeatletter\displaymath@qed
\end{equation*}
%for representatives of heights provided by the Height Machine. 


%% Recall the projection $\rho_1 \colon \IP_k^n \times \IP_k^n \times
%% \IP_k^m \rightarrow \IP_k^n \times \IP_k^m$ defined by $(P,Q,s)
%% \mapsto (P,s)$, which induces an isomorphism $\rho_1|_X \colon X
%% \rightarrow X_N$. 
%% Let $U_l$ be a  Zariski open dense subset of $X$ with $U_l\subset
%% \rho_1(U'_l)$. 
%% %Let $U_l := \rho_1|_X^{-1}(U'_l \cap X_N)$. 
%% Then for any $P \in U_l(\IQbar)$, we have
%% \[
%% h([N]P) \ge c_1 N^2 h(\pi(P)) - c_2(l)
%% \]
%% by \eqref{EqAuxHtIneqAux} in view of \eqref{EqHeightTransfer}!!!!. This is precisely the desired height inequality as $N = 2^l$.
% (so that we can denote by abuse of notation $U_l = U_N$ and $c_2(l) = c_2(N)$).

%%% Local Variables:
%%% TeX-master: "main"
%%% End:
